% !TEX root = ../main.tex
%----------------------------------------------------------
% appendix-alternative-scores.tex — Appendix A
% Moved from §4.4 (IRLS) and §4.5 (Possibilistic Climatology)
%----------------------------------------------------------

\appendix
\section{Alternative Scores and Baselines}\label{sec:appendix_alt}

%----------------------------------------------------------
\subsection{Imprecise Ranked Logarithmic Score}\label{sec:irls}

For $K$ ordered categories, the cumulative possibility and necessity
bounds are:
\begin{align}
U_k &= \max_{i \leq k}\, \pi(i) \\
L_k &= m \cdot \bigl[1 - \max_{i > k}\, \pinorm(i)\bigr]
\end{align}
with $\max$ over an empty set defined as $0$, so that $U_K = m$ and
$L_K = m$.

The \textit{imprecise ranked logarithmic score} is:
\begin{equation}
\mathrm{IRLS} = \frac{1}{K} \sum_{k=1}^{K}
    \bigl[e_k \cdot S(U_k) + (1 - e_k) \cdot S(1 - L_k)\bigr]
\label{eq:irls}
\end{equation}
where $S(p) = -\log_2(\max(p, \epsilon))$ and
$e_k = \mathbf{1}\{c_{\mathrm{obs}} \leq k\}$.  When $L_k = U_k$ for
all $k$, IRLS reduces to the standard ranked logarithmic score.

\textbf{Properties:}
\begin{itemize}[leftmargin=2em]
\item $U_k$ is non-decreasing in $k$ (maximum over a growing set).
\item $L_k$ is non-decreasing in $k$ (complement of a maximum over a
    shrinking set).
\item $L_k \leq U_k$ for all $k$.
\end{itemize}

The IRLS bears the same relationship to ILS that CRPS bears to the Brier score: it integrates the threshold-specific score across all category boundaries, producing a single summary measure that rewards both sharpness and calibration across the full distribution rather than at a single threshold.

\textbf{Note:} IRLS is formally defined here for completeness, but its
detailed application is deferred to future work with a larger
reforecast dataset.

%----------------------------------------------------------
\subsection{Possibilistic Climatology}\label{sec:poss_climatology}

A possibilistic baseline analogous to probabilistic climatology is
needed for skill scores.  The recommended construction derives
\textit{consonant} possibility distributions from the climatological CDF.

For continuous variables with climatological CDF $F_{\mathrm{clim}}(z)$:
\begin{equation}
\pi_{\mathrm{clim}}(z) = 1 - 2\,\bigl|F_{\mathrm{clim}}(z) - 0.5\bigr|
\label{eq:poss_clim}
\end{equation}
This peaks at the climatological median ($\pi_{\mathrm{clim}} = 1$ where
$F_{\mathrm{clim}} = 0.5$) and produces central nested $\alpha$-cut sets
$S_{\mathrm{clim}}(\alpha) = [Q_{\mathrm{clim}}(\alpha/2),\;
Q_{\mathrm{clim}}(1 - \alpha/2)]$.

For $K$ ordered categories with smoothed climatological frequencies
$p_k$, the \textit{categorical mid-rank method} is recommended: compute
cumulative frequency $F_{k-1} = \sum_{j < k} p_j$, category mid-rank
$u_k = F_{k-1} + 0.5\, p_k$, and define:
\[
\pi_{\mathrm{clim}}(k) = 1 - 2\,|u_k - 0.5|
\]

Climatological ignorance is zero by convention
($H_{\Pi,\mathrm{clim}} = 0$): climatology is a reference distribution,
not a live forecasting system.
The consonant construction is the least-committed possibility distribution consistent with a given CDF: its nested $\alpha$-cuts are central intervals that encode no shape preference beyond what the cumulative frequencies imply.  For applications with multimodal climatologies or strongly asymmetric base rates, alternative constructions (e.g., non-nested random sets) may be more appropriate; the choice of climatological reference should be justified on physical grounds.  Any $\ign > 0$ from an operational system therefore reflects genuine model uncertainty beyond this baseline---the system's admitted ignorance is measured relative to a reference that, by convention, claims complete knowledge.

Persistence---the previous day's distribution used as today's
forecast---is another valid baseline.  A system that cannot outperform
persistence has negative skill relative to the simplest possible
dynamical reference.
